\subsection{DRP: conmutación regional, retorno y custodia}\label{sec:sd8-drp}
El par seleccionado es AWS São Paulo, \texttt{sa-east-1}, y AWS México Central, \texttt{mx-central-1}, activo-pasivo preparado según DR-4 de SD4. La separación orientativa es 7.613,84 km entre referencias urbanas, no centros de datos ni latencia. Reduce amenazas físicas locales; no elimina dependencia común de AWS, organización/IAM, DNS, federación, claves, despliegue, corrupción o rutas de enlace. SRE comprueba esas dependencias y Seguridad autorizaciones de tratamiento Brasil/México y custodia; un traslado regional no acredita permiso legal por sí mismo. No se mantiene Google Cloud Santiago como un segundo destino simultáneo ni México como región temporal.

\textbf{DR-01, detección y decisión.} La prueba mide RTO desde la interrupción y RPO desde el último corte común recuperable; SRE abre incidente y comunica al CLIENTE, Seguridad analiza integridad y Proyecto coordina consecuencias. Las alarmas de SD4 usan latido SQL y objeto, manifiestos y observación cada minuto: más de 120 s en dos muestras avisa; más de 300 s en dos genera incidencia; más de 600 s en una escala y suspende liberaciones/cargas no críticas; 900 s, fallo o ausencia de dos muestras válidas bloquean declaración conforme de RPO. Se contempla reloj no fiable y objetos pendientes; \texttt{ReplicaLag} no demuestra por sí solo copia de toda la evidencia. El incumplimiento del RPO se registra sin alterar su mínimo.

\textbf{DR-02, exclusión y corte.} Arquitectura/SRE ejecutan el procedimiento de aislamiento, detienen emisores y escrituras del origen y conservan ID/época de autoridad. Antes de promover, Datos verifica corte común de los dos almacenes, objetos, permisos, revocaciones, eliminaciones y bandejas de hechos. Una partición que impida comprobar exclusión del escritor anterior no autoriza promoción; se mantiene operación local segura y se escala. DNS por sí solo no impide a un escritor antiguo producir efectos. Si la evidencia del diseño no resuelve exclusión, identidad de emergencia o clave, queda una brecha de arquitectura y no se declara alcanzado el RTO por saltar el control.

\textbf{DR-03, preparación y restitución.} SRE promueve réplicas según SD4, ajusta capacidad al perfil productivo vigente, activa infraestructura versionada y comprueba cuotas, certificados, claves, identidad regional y ruta de acceso desde cliente y recinto. No se replica una contraseña ni un secreto externo por suponer que el proveedor lo haga. El pool de contingencia conserva identidad estable y mecanismos de recuperación preparados de SD4; federación caída no se usa como único acceso para respuesta. Seguridad comprueba permisos mínimos y autenticación de responsables de emergencia. Colas se reconstruyen desde inbox/efecto/outbox durables: los servicios de transporte no son la única copia. Se consultan respuestas externas inciertas por ID antes de repetir, y se bloquea un efecto sin idempotencia/consulta hasta conciliación competente.

Calidad y referente funcional prueban recepción/amarra, salida/regreso, presencia/retirada escolar, combustible, avisos, integridad y consulta/efecto administrativo crítico. Se reanudan consumidores y se abre tráfico sólo tras esas comprobaciones. SRE comunica restablecimiento y restricciones restantes; la duración del incidente continúa hasta recuperación efectiva del negocio, no hasta que un proceso arranque. La Contraparte recibe procedimiento y entrenamiento para solicitar, autorizar o coordinar la conmutación desde la consola gestionada; la ejecución técnica pertenece al servicio Synaptix, sin exigir un administrador AWS local inexistente. \parencites[RT-07.01--07]{sd8-bt}[numerales 2.4 y 17.6]{sd8-c07}

\textbf{Presupuesto técnico de RTO, hipótesis por validar.} Se asignan 15 min a detección/decisión, 30 a exclusión y corte común, 30 a promoción, 45 a capacidad/rutas/identidad y 90 a prueba integrada y reconciliación: 210 min; los 30 restantes son margen de ese escenario, no una nueva reserva de calendario contractual. Las fases solapadas no se descuentan sin prueba. Fallo de clave, identidad, datos o exclusión puede exceder esa secuencia: el plan conserva la brecha y corrige antes de producción. No se presenta el presupuesto como medición realizada ni garantía de capacidad faltante.

\textbf{DR-04, retorno.} Se estabiliza el primario sin servir escritura, se realiza copia/replicación desde la autoridad actualmente activa en México, se reconcilian hechos, objetos, permisos y efectos externos, y se ensaya el retorno de los pendientes de borde. Datos verifica corte y Calidad negocio; SRE ejecuta un cambio de autoridad autorizado con aislamiento del escritor anterior y observación reforzada. Se mantiene origen de contingencia protegido hasta confirmar eficacia y aceptación; no se reutiliza una réplica divergente ni se eliminan evidencias para aparentar normalidad. Proyecto registra residual observado, duración, último dato recuperable, condiciones de retorno y acciones correctivas.

\subsection{Respaldo, restauración y conservación}\label{sec:sd8-backup}
Se aplica 3-2-1-1-0: tres copias recuperables del conjunto, dos medios distintos, una fuera de sitio, una inmutable o fuera de línea y cero errores en la verificación. Datos/SRE mantienen producción/réplica, respaldo versionado y custodia externa independiente prevista en SD4, incluyendo original, documentos y diario local pendiente. Una réplica que propaga corrupción no sustituye la copia de recuperación; tres referencias al mismo objeto no son tres copias. La política requiere inventario de medio, ubicación, propietario, corte, huella, cifrado, clave, retención, expiración y resultado de prueba. La variante de custodia y sus recursos se verifican antes de confiarle recuperación, sin afirmar contrato ya firmado.

Claves de respaldo se custodian de forma independiente del conjunto respaldado, con acceso nominal, doble control para recuperación de emergencia y prueba de descifrado desde el entorno de restauración. Copias inmutables tienen regla de retención que impide borrado/modificación aun con credenciales administrativas comprometidas; Seguridad ensaya el intento denegado. Una pérdida de clave, copia incompleta o restauración con datos vencidos impide abrir tráfico. Se restaura aislado y se reaplican revocaciones/eliminaciones autorizadas y bloqueos probatorios antes de servir. \parencites[art.~20]{sd8-ba}[RT-07.09--13]{sd8-bt}

Frecuencia propuesta para los doce dominios de SD5: copia diaria, recuperación a punto en el tiempo de 35 días en nube y doce copias mensuales; subconjunto local diario de treinta días. La tabla de tiempos de SD5 reúne nueve grupos de recuperación: Clientes/contratos, Operación marítima, Escuela, Varadero/contratistas, Combustible, Ambiente/consumos, Finanzas/cobranza, Infraestructura y Auditoría/seguridad. Sus tiempos de restauración individual estimados por SD5 son, en ese orden, 74, 54, 71, 65, 47, 54, 71, 51 y 47 minutos para su volumen histórico provisional. Son escenarios por grupo y no se suman para restauración conjunta, estimada en 233,33 minutos con los supuestos explicados en 8.2. Antes de aceptar RT-07.13, Datos desglosa esa estimación por NAV, AMA, ESC, CON, BOY, VAR, CTR, AMB, FAC, ADM, AUT y ALE, incluyendo avisos y sus acuses. Los nueve grupos no acreditan todavía tiempo completo por cada dominio; se conservan como cálculo parcial y se agregan subconjuntos y procesos nuevos al inventario antes de respaldarlos; los documentos de evidencia mantienen retención propia, que puede superar treinta días, 35 días o doce meses. No se confunde rotación de respaldo con retención legal de fuente.

SRE verifica copia diaria, corte y huella; mensualmente restaura muestra representativa con claves, permisos, documentos y hechos pendientes, mide el tiempo efectivo y Calidad valida integridad. Un error invalida la prueba y abre remediación con responsable/fecha y repetición; cero errores es criterio de aceptación, no un resultado previamente obtenido. Los datos locales no confirmados no se borran por antigüedad y se incluyen en respaldo/replicación local antes de purgar; al reconectar se contrastan manifiestos de destino. Las muestras de menores se minimizan y protegen sin convertir datos de prueba en una divulgación.

\subsection{Pruebas, ventanas y criterio de recuperación}\label{sec:sd8-pruebas-continuidad}
Antes de cada producción E1/E2, Calidad coordina con SRE/Datos/Seguridad pruebas de recuperación regional, dos enlaces WAN retirados durante 24 h, terminales 12 h e Ilque 8 h, sincronización menor o igual a treinta minutos, nodo local primario y alternativa, mensajes perdidos/duplicados, caída de identidad, clave no disponible, copia corrupta, objeto sin réplica y eliminación/restauración. Se mide inicio/fin, RTO/RPO, corte de datos, recuentos y huellas, cola y latencia de negocio. Cada prueba conserva versión, carga representativa, ejecutor, verificador independiente cuando corresponde, permiso, incidencia y decisión. No se presenta su planificación como resultado ejecutado.

En operación se realizan al menos dos conmutaciones regionales reales al año, con informe RTO/RPO y corrección de brechas; resiliencia por instancia, zona, dependencia, latencia y disco al menos semestral; restauración representativa mensual; seguridad ofensiva por tercero independiente anual y antes de cada producción. SRE incorpora las pruebas en 1.12.4 y Calidad en paquetes de aceptación de SD7/SD9, evitando duplicar una ejecución. Se programan fuera de congelamiento y ventanas críticas con aviso de al menos diez días hábiles; la fecha se anticipa para conservar intervalo máximo, no se posterga por conveniencia. Una emergencia no requiere esperar ese aviso de mantenimiento. \parencites[arts.~20 y 21.3]{sd8-ba}[RT-07.07/12; RT-10.05--07; RT-11.20; cap.~20]{sd8-bt}

El CLIENTE valida el modo de operación y comunicación; el servicio gestionado ejecuta técnica. El retorno requiere recuperar flujos y accesos, conciliar hechos alternativos y pendientes, comprobar retenciones y restricciones, retirar doble autoridad y emitir comunicación de restablecimiento. Una condición impeditiva mantiene bloqueo incluso con puntuación residual baja. Los registros de prueba y eficacia alimentan T-16 y los comités; sólo entonces se propone residual observado o reapertura. Permanecen visibles las brechas de capacidad, calendario y diseño señaladas por SD4/SD5/T-15, sin trasladarlas a una aceptación tácita.
